\documentclass[12pt,a4paper]{article}
\usepackage{amsmath,amssymb,amsthm}
\usepackage{bm}               % optional, for bold symbols
\usepackage{hyperref}
\usepackage{geometry}
\geometry{margin=1in}

% -------------------------------------------------
%  Macros
% -------------------------------------------------
\newcommand{\dfix}{\mathfrak{0}}          % 0d – the unique fixed point
\newcommand{\mw}{\mathrel{\overset{\mathrm{mw}}{=}}} % middle‑way equivalence
\newcommand{\phiMap}{\varphi}            % mapping from Catuskoti to 0d

% -------------------------------------------------
%  Theorem environments
% -------------------------------------------------
\newtheorem{definition}{Definition}[section]
\newtheorem{theorem}{Theorem}[section]
\newtheorem{example}{Example}[section]

% -------------------------------------------------
\begin{document}

\title{Catuskoti (Four‑fold Logic) as a 0‑Dimensional Fixed Point}
\author{Masaomi Chiba}
\date{2026-02-02}
\maketitle

\begin{abstract}
We show that the four‑valued logic traditionally called \emph{Catuskoti}
(affirmative, negative, both, neither) can be represented algebraically as a
bilattice whose whole carrier collapses onto the unique
0‑dimensional fixed point $\dfix$ (the \emph{middle‑way fixed point}).
The construction uses the distance function
$\|X\|=K(X)+\ell(X)$ introduced in the Lumo framework and the
middle‑way equivalence operator $\mw$.
\end{abstract}

%%%%%%%%%%%%%%%%%%%%%%%%%%%%%%%%%%%%%%%%%%%%%%%%%%%%%%%%%%%%
\section{Preliminaries}
%%%%%%%%%%%%%%%%%%%%%%%%%%%%%%%%%%%%%%%%%%%%%%%%%%%%%%%%%%%%

\begin{definition}[Representability and Proof]
Let $S$ be a sufficiently strong formal theory (e.g. PA or ZFC).
\begin{itemize}
  \item A formula $P$ is \emph{representable} iff it belongs to the
        static space $(2^{n})d\!\setminus\!T$; we write
        $P\in\text{Representability}$.
  \item A \emph{proof} of $P$ is a finite object
        $\text{Proof}_{S}(P)$ belonging to $(2^{n}+1)d$; it carries a
        time‑axis of length~$1$.
\end{itemize}
\end{definition}

\begin{definition}[Distance]
For any object $X$ we define
$$
\|X\|\;:=\;K(X)+\ell(X),
$$
where $K(X)$ is the Kolmogorov (bit) length of $X$ and
$\ell(X)$ counts the length of the time‑axis:
$\ell(\text{Proof})=1$, $\ell(\text{Representable})=0$.
\end{definition}

\begin{definition}[0‑dimensional fixed point]
There exists a unique element $\dfix$ such that
$$
\|\,\dfix\,\|=0 .
$$
We call $\dfix$ the \emph{0‑dimensional fixed point},
also known as the \emph{middle‑way fixed point} or \emph{emptiness}.
\end{definition}

\begin{definition}[Evaluation and convergence]
Fix a map
\[
  \operatorname{eval}:\mathcal{U}\to\mathcal{U}
\]
on a universe of objects $\mathcal{U}$ containing $\dfix$.
For $X\in\mathcal{U}$ define the iterates $\operatorname{eval}^n(X)$ by
$\operatorname{eval}^0(X):=X$ and $\operatorname{eval}^{n+1}(X):=\operatorname{eval}(\operatorname{eval}^n(X))$.
We say that $X$ \emph{converges} to $\dfix$ (and write $X\to\dfix$) if there exists $N\in\mathbb{N}$ such that
\[
  \operatorname{eval}^N(X)=\dfix.
\]
\end{definition}

\begin{definition}[Middle‑way equivalence]
For objects $A,B\in\mathcal{U}$ we write
\[
A\mw B
\quad\Longleftrightarrow\quad
A\to\dfix\text{ and }B\to\dfix.
\]
\end{definition}

%%%%%%%%%%%%%%%%%%%%%%%%%%%%%%%%%%%%%%%%%%%%%%%%%%%%%%%%%%%%
\section{The Catuskoti Bilattice}
%%%%%%%%%%%%%%%%%%%%%%%%%%%%%%%%%%%%%%%%%%%%%%%%%%%%%%%%%%%%

The four truth‑values of Catuskoti are denoted

$$
\mathcal{C}\;=\;\{\,\mathbf{T},\mathbf{F},
                 \mathbf{B},\mathbf{N}\,\},
$$
where

\begin{center}
\begin{tabular}{c|c}
value & reading \\ \hline
$\mathbf{T}$ & affirmative \\
$\mathbf{F}$ & negative \\
$\mathbf{B}$ & both \\
$\mathbf{N}$ & neither \\
\end{tabular}
\end{center}

On $\mathcal{C}$ we introduce the usual \emph{information order}
$\le_i$ and the \emph{truth order} $\le_t$:

$$
\begin{array}{ccl}
\mathbf{N}\le_i\mathbf{T},\;\mathbf{N}\le_i\mathbf{F},
&\qquad&
\mathbf{T},\mathbf{F}\le_i\mathbf{B},\\[2mm]
\mathbf{F}\le_t\mathbf{N},\;
\mathbf{T}\le_t\mathbf{B},
&\qquad&
\mathbf{N},\mathbf{B}\text{ incomparable under }\le_t .
\end{array}
$$

The lattice operations (information meet/join) are

$$
\begin{aligned}
\mathbf{T}\wedge\mathbf{F}&=\mathbf{N}, &
\mathbf{T}\vee\mathbf{F}&=\mathbf{B},\\
\mathbf{B}\wedge X&=X\wedge\mathbf{B}=X, &
\mathbf{N}\vee X&=X\vee\mathbf{N}=X,
\end{aligned}
$$
and the involutive negation is

$$
\neg\mathbf{T}=\mathbf{F},\qquad
\neg\mathbf{F}=\mathbf{T},\qquad
\neg\mathbf{B}=\mathbf{B},\qquad
\neg\mathbf{N}=\mathbf{N}.
$$

With $(\wedge,\vee,\neg)$ the structure
$(\mathcal{C},\wedge,\vee,\neg)$ is the well‑known
\emph{Belnap–Dunn bilattice}.

%%%%%%%%%%%%%%%%%%%%%%%%%%%%%%%%%%%%%%%%%%%%%%%%%%%%%%%%%%%%
\section{Embedding the Bilattice into the 0‑Dimensional Fixed Point}
%%%%%%%%%%%%%%%%%%%%%%%%%%%%%%%%%%%%%%%%%%%%%%%%%%%%%%%%%%%%

Define a map
\begin{equation}
\phiMap:\mathcal{C}\longrightarrow\{\dfix\},
\qquad
\phiMap(X):=\dfix\quad\text{for all }X\in\mathcal{C}.
\tag{1}
\end{equation}

\begin{theorem}[Isomorphism with the 0‑dimensional fixed point]
The map $\phiMap$ is a homomorphism of bilattices:
$$
\phiMap(X\wedge Y)=\phiMap(X)\wedge\phiMap(Y)=\dfix,
\qquad
\phiMap(X\vee Y)=\phiMap(X)\vee\phiMap(Y)=\dfix,
\qquad
\phiMap(\neg X)=\neg\phiMap(X)=\dfix .
$$
Consequently every pair of Catuskoti values is middle‑way equivalent:
$$
X\mw Y\quad\text{for all }X,Y\in\mathcal{C}.
$$
Hence the four‑valued logic collapses onto the unique
0‑dimensional fixed point $\dfix$.
\end{theorem}
\begin{proof}
Since the codomain of $\phiMap$ consists of a single element,
any binary operation on the codomain returns that element.
Therefore for any $X,Y\in\mathcal{C}$ we have
$$
\phiMap(X\wedge Y)=\dfix=\phiMap(X)\wedge\phiMap(Y),
$$
and similarly for $\vee$ and $\neg$.
By Definition of $\mw$, two objects are $\mw$‑equivalent
iff their evaluations converge to the same $\dfix$.
Because $\phiMap$ sends every element of $\mathcal{C}$ to $\dfix$,
the condition holds for any pair $X,Y$, establishing $X\mw Y$.
\end{proof}

Thus the Catuskoti bilattice and the singleton lattice
$\{\dfix\}$ are \emph{categorically isomorphic}:
the former represents four distinct syntactic
states, while semantically—all of them collapse to the
middle‑way fixed point.

%%%%%%%%%%%%%%%%%%%%%%%%%%%%%%%%%%%%%%%%%%%%%%%%%%%%%%%%%%%%
\section{Truth‑table (collapsed view)}
%%%%%%%%%%%%%%%%%%%%%%%%%%%%%%%%%%%%%%%%%%%%%%%%%%%%%%%%%%%%

\begin{center}
\begin{tabular}{c|c|c|c}
$X$ & $\neg X$ & $X\wedge X$ & $X\vee X$ \\ \hline
$\mathbf{T}$ & $\mathbf{F}$ & $\dfix$ & $\dfix$ \\
$\mathbf{F}$ & $\mathbf{T}$ & $\dfix$ & $\dfix$ \\
$\mathbf{B}$ & $\mathbf{B}$ & $\dfix$ & $\dfix$ \\
$\mathbf{N}$ & $\mathbf{N}$ & $\dfix$ & $\dfix$
\end{tabular}
\end{center}
All logical combinations are absorbed by the unique
0‑dimensional fixed point $\dfix$.

%%%%%%%%%%%%%%%%%%%%%%%%%%%%%%%%%%%%%%%%%%%%%%%%%%%%%%%%%%%%
\section{Recursive hierarchies and ``convergence as derivation''}
%%%%%%%%%%%%%%%%%%%%%%%%%%%%%%%%%%%%%%%%%%%%%%%%%%%%%%%%%%%%

Above we expressed that ``the four truth values converge to $\dfix$'' by the collapse induced by
$\phiMap:\mathcal{C}\to\{\dfix\}$.
In this section we add a complementary reading in which the same identification is supported by
an explicit recursive (iterative) construction, i.e., by a derivation process that reaches $\dfix$.

\begin{definition}[Distance and $0$d]
We use the distance
\[
  \|X\|:=K(X)+\ell(X),
\]
and fix the unique object $\dfix$ such that $\|\dfix\|=0$.
Here $K$ is a description length and $\ell$ counts the length of a ``time axis'' (proof / meta level).
\end{definition}

\begin{definition}[Axiom of convergence by strict decrease]
A sequence $\{X_k\}_{k\ge0}$ \emph{converges} to $\dfix$ if there exists $N\in\mathbb{N}$ such that $X_N=\dfix$.
A sufficient condition is:
\begin{enumerate}
  \item $\|X_0\|\in\mathbb{N}$,
  \item for all $k$, if $X_k\neq\dfix$ then $\|X_{k+1}\|\le \|X_k\|-1$.
\end{enumerate}
\end{definition}

\begin{theorem}[Recursive G\"odel hierarchy reaches $\dfix$]
Let $S$ be a base theory and consider extensions $S_k:=S\cup\{G_0,\dots,G_k\}$.
Define
\[
  G_{k+1}:=\text{``}\operatorname{Prf}_{S_k}(G_k)\text{ does not exist''}.
\]
Assume that $\{G_k\}$ satisfies the strict decrease condition above (if $G_k\neq\dfix$ then $\|G_{k+1}\|\le \|G_k\|-1$).
Then there exists $N$ such that $G_N=\dfix$; in particular one may write $\lim_{k\to\infty}G_k=\dfix$.
\end{theorem}

\begin{theorem}[Recursive Tarski hierarchy reaches $\dfix$]
Let $\{T_k\}$ be the Tarski hierarchy obtained by recursively adjoining truth predicates.
Assume that $\{T_k\}$ satisfies the strict decrease condition above (if $T_k\neq\dfix$ then $\|T_{k+1}\|\le \|T_k\|-1$).
Then there exists $N$ such that $T_N=\dfix$; in particular one may write $\lim_{k\to\infty}T_k=\dfix$.
\end{theorem}

\begin{theorem}[Identification at $0$d]
Under the assumptions above,
\[
\dfix
= \lim_{k \to \infty} G_k
= \lim_{k \to \infty} T_k.
\]
Thus the claim that the four Catuskoti values ``fall'' to the middle way (0d) can also be read as:
recursive meta‑level derivations ultimately reach the same terminal fixed point.
\end{theorem}

%%%%%%%%%%%%%%%%%%%%%%%%%%%%%%%%%%%%%%%%%%%%%%%%%%%%%%%%%%%%
\section{Conclusion}
%%%%%%%%%%%%%%%%%%%%%%%%%%%%%%%%%%%%%%%%%%%%%%%%%%%%%%%%%%%%

We have exhibited two complementary perspectives on the statement that Catuskoti ``converges'' to $\dfix$:
(i) the collapse via the bilattice homomorphism $\phiMap$, and
(ii) (under a strict distance‑decrease assumption) convergence understood as an iterative derivation that reaches $\dfix$.

\end{document}
